\documentclass[11pt,letterpaper]{article}

\usepackage[margin=1in]{geometry}
\usepackage{amsmath}
\usepackage{graphicx}
\usepackage{hyperref}
\usepackage{booktabs}
\usepackage{array}
\usepackage{xcolor}
\usepackage{enumitem}
\usepackage{caption}
\usepackage{subcaption}
\usepackage{natbib}
\usepackage{microtype}
\usepackage{parskip}

\hypersetup{
  colorlinks=true,
  linkcolor=blue!60!black,
  citecolor=blue!60!black,
  urlcolor=blue!60!black,
}

\title{\textbf{Brain Gene Expression Explorer}\\
  \large Interactive 3D Visualization of Spatially-Resolved\\
  Gene Expression Across a Whole-Brain Parcellation}

\author{Thiago Viegas\\
  \small New York University}

\date{\today}

% ─────────────────────────────────────────────────────────────────────────────
\begin{document}

\maketitle

\begin{abstract}
Understanding where genes are expressed across the human brain is fundamental to
neurogenomics research, yet the data are high-dimensional, spatially embedded, and
produced by two complementary but methodologically distinct resources: the
Genotype-Tissue Expression (GTEx) project and the Allen Human Brain Atlas (AHBA).
We present \emph{Brain Gene Expression Explorer}, an interactive 3D visualization
tool that maps gene expression estimates onto a 156-parcel whole-brain atlas and
renders them in a browser-native WebGL viewer.
The application displays four side-by-side panels that simultaneously show
(V1)~raw GTEx measurements at held-out parcels, (V2)~cross-validated
out-of-sample reconstructions, (V3)~dense whole-brain model predictions,
and (V4)~a naive-model AHBA reference baseline,
enabling direct visual comparison of data quality and model performance.
A residual mode switches V2 to a prediction-error map and V3 to a
whole-brain-fit minus AHBA-reference difference map, both rendered on
a diverging colormap with per-parcel breakdown tooltips on click.
Users can explore 88 high-variance genes across five donors under three
model variants (naive, DLAM, PLAM) with real-time camera controls,
a half-brain sagittal clip, and independent opacity sliders for cortex,
subcortex, and cerebellum.
\end{abstract}

% ─────────────────────────────────────────────────────────────────────────────
\section{Introduction \& Motivation}
\label{sec:intro}

\subsection{Problem Definition}

Mapping gene expression to brain anatomy is a core challenge in neuroinformatics.
Two major resources address this:
the \textbf{Allen Human Brain Atlas (AHBA)}~\citep{hawrylycz2012} provides
genome-wide mRNA expression for roughly 900 cortical and subcortical sites
from a small number of donors (typically 6), while the
\textbf{Genotype-Tissue Expression (GTEx)} project~\citep{gtex2013}
provides bulk RNA-seq from hundreds of brain samples but with only coarse
anatomical labelling.
Computational pipelines that integrate and spatially impute these datasets produce
per-parcel expression estimates for each gene across the whole brain;
however, they output dense numerical arrays ($\sim$150 parcels $\times$ 13{,}618 genes)
that are difficult to interpret without purpose-built visual tooling.

\subsection{Why It Matters}

Gene expression shapes brain function, development, and disease.
Spatially resolved expression maps are used to nominate candidate genes for
region-specific phenotypes, validate cell-type deconvolution models, and
generate hypotheses about transcriptomic gradients.
Existing neuroimaging viewers (FSLeyes, Connectome Workbench GUI) handle
continuous volumetric or surface data but do not natively support
switching between hundreds of gene ``channels'' or comparing model-variant
outputs side by side.
A fast, interactive viewer tailored to the parcel-level granularity of these
pipelines accelerates the iterative analysis loop for researchers and students
who are not VTK specialists.

\subsection{Target Users / Use Case}

The primary users are neuroinformatics researchers and graduate students working
with the GTEx--AHBA integration pipeline at NYU.
Typical use cases include:
\begin{enumerate}[nosep]
  \item inspecting whether a model's predictions agree with sparse measured values
        for a specific gene and donor;
  \item comparing two model architectures (e.g.\ DLAM vs.\ PLAM) at a glance;
  \item generating figures for presentations by exporting specific views.
\end{enumerate}

% ─────────────────────────────────────────────────────────────────────────────
\section{Related Work}
\label{sec:related}

\subsection{Neuroimaging Viewers}

Several established tools render brain volumes and surfaces interactively.
\textbf{FSLeyes}~\citep{mccarthy2021} is a desktop application for 4D
NIfTI data that supports multiple overlays and time-series playback.
\textbf{Connectome Workbench}~\citep{marcus2011} handles surface-based
fMRI data on standard meshes (fsLR) but requires the user to load one
metric file per gene and per model variant—an impractical workflow for
13{,}618 genes.
\textbf{BrainPainter}~\citep{marinescu2019} produces static renderings from
label maps, with no interactive exploration.

\subsection{Gene Expression Atlases}

The \textbf{Allen Brain Atlas (ABA) online viewer}
provides a reference browser for the original AHBA in-situ hybridization data
at voxel resolution but does not support RNA-seq estimates, parcel-level
aggregations, or model output comparison.
\textbf{abagen}~\citep{markello2021} is a Python toolkit for processing AHBA data;
it outputs DataFrames rather than interactive visualizations.
\textbf{Nilearn}~\citep{abraham2014} offers glass-brain and surface plots
in static matplotlib figures, useful for publications but not for interactive
exploration.

\subsection{Web-based 3D Viewers}

\textbf{Plotly Dash / Plotly 3D surface} tools can render brain surfaces
but struggle with the $\sim$100k-vertex meshes used by standard cortical
reconstructions without significant LOD engineering.
Trame (used in this project) renders via VTK's server-side pipeline and
streams only compressed diff frames to the client, making it suitable for
dense mesh workloads on a local server.

\subsection{Gap Addressed}

No existing open tool combines: (a) parcel-level gene expression from
RNA-seq integration pipelines, (b) multi-panel model comparison,
(c) fast in-memory gene switching, and (d) a browser-accessible interface
requiring no desktop installation from the end user.
Our application fills this gap for the specific workflow of GTEx--AHBA
spatial imputation.

% ─────────────────────────────────────────────────────────────────────────────
\section{Data Description}
\label{sec:data}

\subsection{Datasets}

\paragraph{GTEx brain RNA-seq.}
Bulk RNA-seq count matrices for five post-mortem brain donors
(GTEX-1117F, GTEX-13OW8, GTEX-11DZ1, GTEX-1B996, GTEX-1JMPZ)
downloaded from the GTEx portal.
Each donor contributes between 8 and 12 anatomically distinct tissue samples,
each mapped to a parcel in the 4S156Parcels atlas.

\paragraph{Allen Human Brain Atlas (AHBA).}
Microarray expression data for~$\sim$900 tissue samples from 6 donors.
The upstream pipeline uses \texttt{abagen} to process and normalise these
data before spatial imputation.

\paragraph{4S156Parcels atlas.}
A whole-brain parcellation with 100 cortical and 56 subcortical regions
defined in MNI152NLin6Asym space~\citep{TODO-atlas-citation}.
Atlas metadata (parcel labels, MNI centroids, structure assignments)
are stored in a reformatted CSV.
The 3D NIfTI volume is used to extract subcortical meshes; the
fsLR32k cortical surface files are used for cortical colouring.

\subsection{Preprocessing}

The upstream integration pipeline produces per-subject \texttt{.npz} archives
(one per model variant) that contain the following arrays, all aligned to
the canonical 150-parcel row order defined in \texttt{parcel\_order.json}
(alphabetical AHBA labels):

\begin{center}
\begin{tabular}{llll}
\toprule
Array key & Shape & Panel & Model-dep.\\
\midrule
\texttt{loro\_truth\_subject\_raw} & (150, $G$) & V1 & No \\
\texttt{loro\_eval\_mask}         & (150,)     & V1 mask & No \\
\texttt{loro\_fused\_subject\_h}  & (150, $G$) & V2 & Yes \\
\texttt{gtex\_mask}               & (150,)     & V2 mask & No \\
\texttt{fullfit\_subject\_h}      & (150, $G$) & V3 & Yes \\
\texttt{imputed\_mask}            & (150,)     & V3 mask & No \\
\texttt{gene\_names}              & (13618,)   & all & No \\
\bottomrule
\end{tabular}
\end{center}

\noindent $G = 13{,}618$ genes.
The viewer further filters to 88 high-variance genes listed in
\texttt{ahba\_100hvg.txt} to keep the dropdown usable.

\subsection{Challenges}

\paragraph{Atlas--data alignment.}
The 4S156Parcels atlas has 156 parcels; 6 lack AHBA microarray coverage
(\texttt{Cerebellar\_Region9}, \texttt{LH-MN}, \texttt{RH-EXA},
\texttt{RH-STH}, \texttt{RH-VeP}, \texttt{RH\_SomMot\_2}) and are
excluded from the \texttt{.npz}.
The application must align the 150-row data arrays to the 156-entry
atlas metadata without relying on integer index order.

\paragraph{Hemisphere convention.}
The upstream pipeline forces all sample MNI $x$-coordinates to be negative
(left hemisphere) before fitting.
As a result, some GTEx parcels are labelled as right-hemisphere structures
in the atlas even though the model treated them as left.
V1 and V2 apply a nearest-neighbour RH$\to$LH remap at render time;
V3 mirrors LH predictions onto their RH counterparts to produce a symmetric
bilateral view.

\paragraph{Mesh scale.}
Subcortical structures range from large nuclei (putamen, thalamus) to tiny
parcels that cannot form a closed mesh at 1\,mm isotropic resolution.
\texttt{LH-MN} and \texttt{RH-HN} were dropped from the mesh cache for this
reason; neither carries GTEx data.

% ─────────────────────────────────────────────────────────────────────────────
\section{Design \& Methodology}
\label{sec:design}

\subsection{Visualization Choices}

\paragraph{Four-panel side-by-side layout.}
The central design decision is to show all four data views simultaneously
on a shared camera.
This uses the \emph{juxtaposition} multiple-views strategy
\citep{gleicher2011}:
V1 shows sparse ground-truth measurements, V2 shows out-of-sample predictions
at those same locations, V3 shows dense whole-brain coverage, and V4 shows
the naive-model AHBA reference used as a fixed baseline.
Placing them side-by-side makes model accuracy and residual patterns visible
at a glance—a user can immediately compare V1 and V2, or V3 and V4.

\paragraph{3D surface encoding.}
We use a \emph{spatial position} primary encoding (parcels rendered in
their anatomically correct MNI locations) plus a \emph{colour} secondary
encoding for expression level.
Position is the most accurate visual channel for spatial anatomy; colour
is the most effective channel for a continuous quantitative variable when
position is already consumed \citep{mackinlay1986}.

\paragraph{Colour.}
Two colormaps are used depending on context:
\begin{itemize}[nosep]
  \item \texttt{viridis} — sequential, for absolute expression values (V1,
        V2 Reconstruction, V3 Whole-brain Fit, V4 AHBA Reference).
        Perceptually uniform~\citep{liu2018}, greyscale-safe, and
        deuteranopia-safe.
  \item \texttt{RdBu\_r} — diverging, for residual maps (V2 Residual and
        V3 Residual).
        Centred at zero so that over-prediction (red) and under-prediction
        (blue) are immediately distinguishable.
\end{itemize}
The colormap range is auto-scaled \emph{per panel} by default; a Shared
mode locks all panels to a joint range, and a WBF mode locks all panels
to V3's range for direct cross-panel comparison.
Unmapped parcels (outside the panel's mask or NaN) are rendered in light
grey ($r{=}0.80, g{=}0.80, b{=}0.80$) so they are distinguishable
from low-expression values without distracting from the data.

\paragraph{Interaction design.}
Seven interaction types are provided:
\begin{enumerate}[nosep]
  \item \emph{Navigation} — mouse drag to orbit, scroll to zoom (all four
        panels simultaneously via a shared VTK camera object).
  \item \emph{Selection} — gene, subject, and model dropdowns change the
        scalar array sliced from the in-memory NPZ.
  \item \emph{Filtering} — the half-brain toggle clips the scene at $x=0$
        and repositions the camera to a lateral view, exposing the
        mesial surface and deep structures.
  \item \emph{Residual mode} — the V2 mode toggle switches V2 from absolute
        reconstruction to a prediction-error map (LORO prediction minus
        harmonized ground truth) and simultaneously switches V3 to a
        whole-brain residual map (WBF minus naive AHBA reference) for
        dlam/plam models; naive stays in absolute mode since it is the
        reference itself.
  \item \emph{Click tooltip} — clicking any parcel displays a floating
        tooltip at the bottom centre of the viewport.  In residual mode the
        tooltip expands to three rows: Truth / Predicted / $\Delta$ for V2,
        and AHBA Reference / WBF / $\Delta$ for V3.
  \item \emph{Opacity adjustment} — three independent sliders control
        cortex (0--1), subcortex (0--1), and cerebellum (0--1) opacity,
        letting the user reveal or hide each compartment without discarding
        the others.
  \item \emph{Colour scale} — a three-way toggle (Local / Shared / WBF)
        controls whether panels auto-scale independently, share a joint
        range, or are all locked to V3's range.
\end{enumerate}

\subsection{Justification Using Course Concepts}

\paragraph{Perception \& pre-attentive features.}
Colour hue is a pre-attentive feature~\citep{ware2004}.
Using a sequential rather than a categorical palette exploits the
ordered colour channel to encode an ordered quantity without requiring
focal attention.
Region boundaries are implied by the parcel mesh edges; no explicit
contour is needed because the visual system groups by proximity and
continuity.

\paragraph{Encoding hierarchy.}
Following Munzner's nested model~\citep{munzner2014}, the task
(compare spatial expression across brain regions and models)
drove the choice of idiom (3D anatomical surface + colour scalar),
which in turn determined the interaction vocabulary (orbit, clip,
channel-switch).

\paragraph{Colour perception.}
The grey NaN colour is chosen to sit near the mean luminance of the
viridis ramp, avoiding the salience bias that would arise if unmapped
parcels were rendered in a high-contrast colour (black, white, or red).

\subsection{Alternative Designs Considered}

\paragraph{2D flatmap.}
A cortical flatmap projection (e.g.\ Mollweide) would eliminate
occlusion and allow all parcels to be visible simultaneously.
We rejected this because subcortical structures cannot be projected
onto a surface flatmap, and the clinical audience is more familiar
with 3D anatomy.

\paragraph{Heatmap matrix.}
A parcels $\times$ genes heatmap would support comparison of many
genes at once but would require the user to maintain a mental map
from parcel label to anatomical location.
The 3D brain surface leverages existing spatial knowledge and reduces
the cognitive load of label-to-anatomy lookup.

\paragraph{Plotly Dash.}
An earlier prototype used Plotly's \texttt{isosurface} trace.
It was discarded because: (a) loading a new gene required re-serialising
the full mesh to JSON on every update, making gene switching slow, and
(b) parcel-level colour mapping was fragile with the isosurface primitive.

% ─────────────────────────────────────────────────────────────────────────────
\section{System Description / Implementation}
\label{sec:system}

\subsection{Tools and Libraries}

\begin{center}
\begin{tabular}{lll}
\toprule
Library & Version & Role\\
\midrule
Python      & 3.10+  & application language \\
PyVista     & latest & VTK Python wrapper; mesh loading, scalar mapping \\
trame       & latest & Python-to-WebGL bridge; event loop and state management \\
trame-vtk   & latest & VTK render serialisation for trame \\
trame-vuetify & latest & Vuetify 3 / Vue 3 UI components defined in Python \\
wslink      & latest & WebSocket server backing trame \\
yabplot     & custom & fsLR32k surface loading and LUT parsing utilities \\
NumPy       & latest & array manipulation \\
matplotlib  & latest & colormap lookup table generation \\
\bottomrule
\end{tabular}
\end{center}

\subsection{Architecture Overview}

The application follows a \emph{server-side rendering} architecture:
all VTK geometry and scalar computations run on the Python server;
the browser client receives compressed differential frame updates
via WebSocket and displays them in a WebGL canvas.
This avoids the need to serialise large mesh data to JSON on every
interaction and keeps startup time low.

\begin{figure}[h]
  \centering
  % \includegraphics[width=0.8\textwidth]{figures/architecture.pdf}
  \fbox{\parbox{0.75\textwidth}{\centering\vspace{1cm}
    [TODO: architecture diagram — Python server $\leftrightarrow$ WebSocket $\leftrightarrow$ Browser WebGL]
  \vspace{1cm}}}
  \caption{Application architecture. All rendering runs in Python/VTK;
           the browser displays a streaming frame buffer.}
  \label{fig:arch}
\end{figure}

\subsection{Key Features}

\paragraph{Efficient gene switching.}
All model variants for the active subject are loaded into memory at
startup (\texttt{\_all\_npz}).
Changing the gene dropdown calls \texttt{np.where(gene\_names == gene)}
to retrieve the column index and slices the pre-loaded arrays in-place.
No disk I/O occurs during interaction, making gene switching instantaneous.

\paragraph{Shared camera.}
A single VTK \texttt{vtkCamera} object is shared across all four
plotter viewports.
Any mouse interaction—orbit, zoom, pan—applies to all panels
simultaneously, preserving spatial registration between views.
Programmatic camera changes (e.g.\ the half-brain lateral snap) are
applied via the VTK camera API directly, bypassing trame's state
serialisation path which can restore a stale camera position on
\texttt{view\_update()}.

\paragraph{Atlas alignment.}
\texttt{data\_loader.load\_alignment()} reads the atlas CSV and
\texttt{parcel\_order.json} and returns a DataFrame whose row order
matches the \texttt{.npz} arrays.
All downstream look-ups index by parcel label string, not by integer
position, making the alignment robust to reordering in either file.

\paragraph{RH$\to$LH remap.}
\texttt{build\_rh\_to\_lh\_map()} pre-computes a 69-entry
nearest-neighbour map from right-hemisphere parcel centroids to
their closest same-structure left-hemisphere counterparts.
V1 and V2 apply this map so that all data appear on the left
hemisphere, consistent with the upstream pipeline's spatial convention.

% ─────────────────────────────────────────────────────────────────────────────
\section{Results \& Insights}
\label{sec:results}

\subsection{Model Comparison}

The four-panel layout revealed systematic differences between
model variants and between predicted and measured values:

\paragraph{LORO cross-validation (V1 vs.\ V2).}
For most genes and donors, V2 (LORO reconstructions) qualitatively
matches V1 (raw GTEx measurements) at the held-out parcels,
confirming that the imputation models generalise to unseen locations.
Cerebellar parcels show the largest deviations (e.g.\ \texttt{DPM1} in
donor GTEX-13OW8: \texttt{Cerebellar\_Region1} V2 = 5.38, V3 = 5.09,
$\Delta = 0.30$), consistent with the known difficulty of imputing
expression in the cerebellum from cortical AHBA donors.

\paragraph{Model variant effects (V3).}
Switching the Model dropdown from \texttt{naive} to \texttt{dlam}
or \texttt{plam} in V3 reveals that the naive baseline produces
spatially inconsistent predictions in the right hemisphere
(visible as noise-like colour variation), while the DLAM and PLAM
models produce smoother, more anatomically plausible bilateral
patterns.
This observation aligns with the design of the models:
DLAM and PLAM incorporate spatial regularisation terms absent from
the naive baseline.

\paragraph{Subcortical expression gradients.}
The cortex opacity slider (lowered to $\sim$0.3) exposes subcortical
structures.
Several genes show strong subcortical--cortical contrasts:
% [TODO: fill in a specific gene and observation from your analysis]
high expression in striatal parcels relative to cortex for \texttt{[GENE]},
and an inverse pattern in thalamic parcels for \texttt{[GENE]}.
These gradients are invisible in the 2D heatmap prototype and became
apparent only in the 3D anatomical view.

\subsection{Residual Maps}

The residual mode (V2 Residual / V3 Residual) proved useful for
diagnosing model-specific biases.
In V2, the prediction-error map ($\texttt{loro\_fused} - \texttt{loro\_truth}$)
highlights which held-out parcels a model over- or under-predicts;
cerebellar parcels show the largest signed errors for the naive baseline.
In V3, the WBF-minus-AHBA-Reference map isolates the spatial pattern of
improvement that DLAM and PLAM achieve over the naive model:
positive (red) regions indicate upward revision relative to the AHBA
reference, negative (blue) regions indicate downward revision.
Because the naive model is the V4 reference, activating residual mode on
the naive subject correctly falls back to absolute display—there is no
residual to compute.

\subsection{Usability Observations}

The half-brain clip proved particularly useful for inspecting the
medial wall and cingulate cortex, which are occluded in the default
lateral view.
The three-way colour scale control (Local / Shared / WBF) addresses the
cross-panel comparison limitation noted in the previous prototype:
Shared mode locks all panels to a joint range so identical colours
represent identical expression values across panels, while Local mode
retains per-panel auto-scaling for maximum within-panel contrast.

% ─────────────────────────────────────────────────────────────────────────────
\section{Limitations \& Future Work}
\label{sec:limits}

\subsection{Current Limitations}

\paragraph{Cortical atlas rendering.}
The cortical surface currently renders as a solid grey shell because
the optional \texttt{wb\_command} pipeline step (Connectome Workbench)
was not completed in time for this submission.
Parcel-level cortical colour mapping requires projecting the volumetric
NIfTI atlas onto the fsLR32k surface; without this step, all cortical
expression is invisible.
Only the 54 subcortical VTK meshes carry parcel colours at this stage.

\paragraph{Per-panel colour scales in Local mode.}
Each panel auto-scales its colormap to the range of the currently
visible data when the Local scale mode is active.
This maximises within-panel contrast but means that identical colours
do not represent identical expression values across panels.
The Shared and WBF scale modes address this for V2/V3/V4 comparison,
but V1's sparse data (9 parcels) still benefits from independent scaling.

\paragraph{No statistical overlay.}
The viewer shows point estimates only.
Uncertainty bounds (e.g.\ bootstrap confidence intervals on model
predictions) are available in the upstream pipeline output but are not
yet exposed in the UI.

\paragraph{Subject re-load latency.}
Switching subjects triggers a full NPZ reload from disk ($\sim$1--2 s
per model variant, 3 variants loaded in parallel at startup).
For the current five-subject corpus this is acceptable, but it would
become a bottleneck with dozens of subjects.

\subsection{Future Work}

\begin{enumerate}[nosep]
  \item \textbf{Complete cortical atlas:} install Connectome Workbench
        and run \texttt{build\_atlas.py} to generate the fsLR32k parcel
        colour map, enabling full cortical visualisation.
  \item \textbf{Uncertainty visualisation:} render confidence intervals
        as a secondary scalar (e.g.\ hatching or secondary opacity).
  \item \textbf{Gene search \& sorting:} replace the dropdown with a
        searchable list sorted by expression variance or by
        differential expression across brain regions.
  \item \textbf{Export:} add a screenshot button that captures the
        current view to a publication-quality PNG with a colour bar legend.
  \item \textbf{Multi-subject overlay:} show average expression across
        all available donors as an additional data view alongside individual subjects.
  \item \textbf{Statistical overlay:} expose bootstrap confidence intervals
        from the upstream pipeline as a secondary scalar channel (e.g.\
        secondary opacity or stippling) to convey prediction uncertainty.
\end{enumerate}

% ─────────────────────────────────────────────────────────────────────────────
\section{Conclusion}
\label{sec:conclusion}

We presented the \emph{Brain Gene Expression Explorer}, an interactive
3D visualization tool that addresses the specific needs of researchers
working with spatially imputed gene expression data from the GTEx--AHBA
integration pipeline.
The key contributions are:

\begin{enumerate}[nosep]
  \item A four-panel side-by-side layout that makes model validation,
        cross-model comparison, and residual analysis visually immediate.
  \item A residual mode that switches V2 to a prediction-error map and V3
        to a whole-brain-fit-minus-AHBA-reference difference map, with
        per-parcel breakdown tooltips on click, supporting rapid diagnosis
        of model-specific spatial biases.
  \item An efficient architecture (server-side VTK, in-memory gene switching,
        shared camera) that keeps all interactions at interactive frame rates.
  \item A principled visual encoding (perceptually uniform sequential and
        diverging colourmaps, pre-attentive spatial mapping, NaN suppression)
        grounded in visualization theory.
  \item A hemisphere-remapping strategy that correctly handles the
        left-hemisphere convention of the upstream pipeline.
\end{enumerate}

The tool is actively used in the NYU Neuroinformatics Lab to diagnose
model behaviour and inspect expression patterns across brain regions,
demonstrating that targeted, domain-specific visualization tooling can
provide scientific value that general-purpose neuroimaging viewers
cannot easily replicate.

% ─────────────────────────────────────────────────────────────────────────────
\bibliographystyle{plainnat}
\begin{thebibliography}{99}

\bibitem[Abraham et al.(2014)]{abraham2014}
Abraham, A., Pedregosa, F., Eickenberg, M., Gervais, P., Mueller, A.,
Kossaifi, J., Gramfort, A., Thirion, B., \& Varoquaux, G. (2014).
Machine learning for neuroimaging with scikit-learn.
\textit{Frontiers in Neuroinformatics}, 8, 14.

\bibitem[GTEx Consortium(2013)]{gtex2013}
GTEx Consortium (2013).
The Genotype-Tissue Expression (GTEx) project.
\textit{Nature Genetics}, 45(6), 580--585.

\bibitem[Gleicher et al.(2011)]{gleicher2011}
Gleicher, M., Albers, D., Walker, R., Jusufi, I., Hansen, C. D., \&
Roberts, J. C. (2011).
Visual comparison for information visualization.
\textit{Information Visualization}, 10(4), 289--309.

\bibitem[Hawrylycz et al.(2012)]{hawrylycz2012}
Hawrylycz, M. J., Lein, E. S., Guillozet-Bongaarts, A. L., et al. (2012).
An anatomically comprehensive atlas of the adult human brain transcriptome.
\textit{Nature}, 489(7416), 391--399.

\bibitem[Liu et al.(2018)]{liu2018}
Liu, Y., Heer, J., \& others (2018).
Somewhere over the rainbow: An empirical assessment of quantitative
colormaps. \textit{IEEE Transactions on Visualization and Computer
Graphics}, 24(1), 330--339.

\bibitem[Mackinlay(1986)]{mackinlay1986}
Mackinlay, J. (1986).
Automating the design of graphical presentations of relational information.
\textit{ACM Transactions on Graphics}, 5(2), 110--141.

\bibitem[Marcus et al.(2011)]{marcus2011}
Marcus, D. S., Harwell, J., Olsen, T., Hodge, M., Glasser, M. F.,
Prior, F., Jenkinson, M., Laumann, T., Curtiss, S. W., \& Van Essen, D. C.
(2011). Informatics and data mining tools and strategies for the
Human Connectome Project.
\textit{Frontiers in Neuroinformatics}, 5, 4.

\bibitem[Marinescu et al.(2019)]{marinescu2019}
Marinescu, R. V., Oxtoby, N. P., Young, A. L., et al. (2019).
BrainPainter: A software for the visualisation of brain structures,
biomarkers and associated pathological processes.
\textit{arXiv preprint} arXiv:1905.08627.

\bibitem[Markello et al.(2021)]{markello2021}
Markello, R. D., Hansen, J. Y., Liu, Z.-Q., Bazinet, V., Shafiei, G.,
Suárez, L. E., Blostein, N., Seidlitz, J., Baillet, S., Satterthwaite,
T. D., \ldots \& Misic, B. (2021).
\textit{neuromaps}: structural and functional interpretation of brain maps.
\textit{Nature Methods}, 19, 1472--1479.

\bibitem[McCarthy(2021)]{mccarthy2021}
McCarthy, P. (2021).
FSLeyes.
\textit{Zenodo}. \url{https://doi.org/10.5281/zenodo.5607715}

\bibitem[Munzner(2014)]{munzner2014}
Munzner, T. (2014).
\textit{Visualization Analysis and Design}.
CRC Press.

\bibitem[Ware(2004)]{ware2004}
Ware, C. (2004).
\textit{Information Visualization: Perception for Design} (2nd ed.).
Morgan Kaufmann.

\end{thebibliography}

\end{document}
